\section{Simple multicurve coefficients of a twisted trace formula}
\label{sec:coefficientdistributions}

This section develops a spectral version of the crossing-resolution
idea. We work throughout with a marked, closed, oriented hyperbolic
surface $X$ of genus $g\geq2$ and curvature $-1$. All closed curves
are unoriented, unless conjugacy classes are explicitly being counted
with orientation. A multicurve may have parallel components; its
length and number of components count those multiplicities. The
empty multicurve is included. Cusps and boundary require additional
trace terms and are not included in the assertions below.

The inputs are the twisted Selberg formula, the simple-multicurve
basis of the character ring, and Dylan Thurston's sign-coherence
theorem for the bracelet basis. We prove the intervening length and
coefficient estimates explicitly. The resulting construction is a
deduction from these established inputs, with no novelty claim.

\subsection{Broken geodesics and the effect of reconnection}

Colin de Verdi\`ere's construction approximates loop space by broken
geodesics and applies stationary phase to a complex heat
parametrix~\cite{CdVPeriodic}. For an equal time subdivision the
finite-dimensional energy is
\[
 E_N(x_0,\ldots,x_{N-1})
 =N\sum_{i=0}^{N-1}d(x_i,x_{i+1})^2,\qquad x_N=x_0,
\]
with consecutive points within the injectivity radius. The critical
equations match incoming and outgoing velocities; the critical
value for a constant-speed closed geodesic is $\ell_\gamma^2$.
On a hyperbolic surface every closed geodesic, including a
self-intersecting one, has positive transverse length Hessian.
For a unit-speed parametrization and a periodic normal variation
$f(s)N(s)$ this is
\[
 \delta^2L_\gamma(f,f)
 =\int_0^{\ell_\gamma}\bigl((f')^2+f^2\bigr)\,\dd s.
\]
Consequently simplicity cannot be selected by retaining only
positive Hessians in the ordinary loop space.

Reconnection enlarges that space. In a small embedded disk of radius
$r$ about a transverse crossing of angle $\theta\in(0,\pi)$, the two
original diameters have total length $4r$. The two noncrossing
pairings have lengths
\begin{align}
 L_0&=4\arsinh\bigl(\sinh r\sin(\theta/2)\bigr),\notag\\
 L_\infty&=4\arsinh\bigl(\sinh r\cos(\theta/2)\bigr).
 \label{eq:coeffuncrossing}
\end{align}
Both are strictly less than $4r$. In particular,
\[
 4r-L_0=4r\bigl(1-\sin(\theta/2)\bigr)+O(r^3).
\]
The descent is first order in the surgery radius. Tightening the
resolved components to geodesics, and discarding contractible
components, cannot increase total length. This is a change of
connectivity, not a negative direction of the original Hessian.
An unconstrained global minimization would select the empty cycle;
we instead keep the different multicurves as algebraic observables.

\subsection{The finite coefficient distributions}

Let $\mathscr C(S)$ be the set of essential simple multicurves on
the underlying marked surface. For
$\rho:\pi_1(S)\to\mathrm{SU}(2)$ set
\begin{equation}
 T_C(\rho)=\prod_{\alpha\subset C}\tr\rho(\alpha),
 \qquad T_\varnothing=1.
 \label{eq:coeffbasis}
\end{equation}
The phrase ``unnegated trace basis'' refers to this sign convention,
not to positivity of the function on $\mathrm{SU}(2)$. The functions
$T_C$ are linearly independent on the unitary character variety and
form a basis of the complex character ring after complexification;
see Charles--March\'e~\cite[Theorem 1.1]{CharlesMarche}. The usual Kauffman
convention uses $-\tr$ per component, so differs by a basis sign.
The relation
\[
 \tr A\,\tr B=\tr(AB)+\tr(AB^{-1})
\]
resolves crossings. For every nontrivial closed curve $\gamma$ there
is a unique finite expansion
\begin{equation}
 \tr\rho(\gamma)=\sum_{C\in\mathscr C(S)}a_C(\gamma)T_C(\rho),
 \qquad a_C(\gamma)\in\ZZ.
 \label{eq:coeffcurveexpansion}
\end{equation}
These integers depend only on topology, not on $X$.

Let $\Delta_\rho$ be the rank-two flat unitary bundle Laplacian,
and write $\lambda_j(\rho)=1/4+r_j(\rho)^2$. For an even
$\varphi\in C_c^\infty(\RR)$ put
\[
 h(r)=\int_\RR\varphi(L)e^{irL}\,\dd L,
\]
and define
\begin{align}
 \mathscr T^X_\varphi(\rho)
 &=\sum_j h(r_j(\rho))
 -\frac{2\operatorname{Area}(X)}{4\pi}
   \int_\RR r h(r)\tanh(\pi r)\,\dd r.\notag
\end{align}
With the two orientations grouped, the twisted Selberg formula is
\begin{equation}
 \mathscr T^X_\varphi(\rho)
 =\sum_{\delta\in\mathcal P_{\rm un}}
   \sum_{k\geq1}
   \frac{\ell_\delta}{\sinh(k\ell_\delta/2)}
   \varphi(k\ell_\delta)\tr\rho(\delta^k).
 \label{eq:coefftwistedSelberg}
\end{equation}
Here $\mathcal P_{\rm un}$ consists of unoriented primitive closed
geodesics. Naud's statement~\cite[Theorem 2.2]{NaudTwisted} sums
over oriented primitive conjugacy classes with denominator
$2\sinh$; grouping inverses gives~\eqref{eq:coefftwistedSelberg}.
The compact length support makes the geometric sum finite.

\begin{definition}[Coefficient distribution]
For $C\in\mathscr C(S)$ define the locally finite signed measure
on $(0,\infty)$
\begin{equation}
 \boxed{
 \sigma_C^X
 =\sum_{\delta\in\mathcal P_{\rm un}}\sum_{k\geq1}
   \frac{\ell_\delta a_C(\delta^k)}
        {\sinh(k\ell_\delta/2)}\,\delta_{k\ell_\delta}.
 }
 \label{eq:coeffdistribution}
\end{equation}
Thus
\[
 [T_C]\mathscr T^X_\varphi=\langle\sigma_C^X,\varphi\rangle.
\]
The notation $\delta_L$ denotes the unit point mass at $L$.
\end{definition}

Local finiteness follows from the finiteness of the closed-geodesic
length set below any bound. The identity term in the trace formula
is a constant function of $\rho$, hence affects only the empty
multicurve coefficient. The input is the \emph{labelled family}
$\rho\mapsto\operatorname{Spec}\Delta_\rho$ on the marked character
variety. A scalar spectrum alone does not retain these coefficients.

\subsection{Length triangularity and the first atom}

\begin{lemma}[Length ordering of resolutions]
\label{lem:coefftriangularity}
If $\delta$ is primitive and $k\geq1$, then
\[
 a_C(\delta^k)\neq0\quad\Longrightarrow\quad
 \ell_X(C)\leq k\ell_X(\delta).
\]
Equality can occur only when $\delta$ is simple and $C=k\delta$,
the multicurve consisting of $k$ parallel copies. Its coefficient
in that case is $1$. If $\delta$ is nonsimple, the inequality is strict.
\end{lemma}

\begin{proof}
For a primitive nonsimple geodesic, completely resolve its crossings.
Each state has a representative obtained by reconnecting the original
arcs in small disks. Formula~\eqref{eq:coeffuncrossing} gives a strict
total saving; tightening and removing contractible loops cannot
increase length. At a multiple intersection, resolve the endpoints
in a small convex disk by a noncrossing pairing. The original pairing
uses intersecting diameters, so at least one new chord is shorter
than a diameter. The same strict conclusion holds. Coefficient
cancellation cannot introduce a state not produced by a resolution.

The product version is also useful: every multicurve in the
expansion of $T_{C_1}\cdots T_{C_j}$ has length at most
$\sum_i\ell_X(C_i)$. Superimpose representatives, separate parallel
components, resolve crossings, and tighten.

For powers use the Chebyshev polynomials
\[
 P_0(x)=2,\quad P_1(x)=x,\quad
 P_{k+1}(x)=xP_k(x)-P_{k-1}(x),
 \qquad \tr(A^k)=P_k(\tr A).
\]
If $\delta$ is simple, $P_k(T_\delta)$ contains only parallel copies
$m\delta$ with $m\leq k$, and its leading term is $T_\delta^k$.
If $\delta$ is nonsimple, every basis term in $\tr\rho(\delta)$
has length strictly below $\ell_\delta$. Apply the product inequality
to each monomial in $P_k$: every nonconstant output has length
strictly below $k\ell_\delta$; the constant term has length zero.
\end{proof}

\begin{proposition}[First atom for a simple support]
\label{prop:coefffirstatom}
For a simple primitive curve $\alpha$ and $m\geq1$,
\begin{equation}
 \boxed{
 \sigma_{m\alpha}^X
 =\frac{\ell_\alpha}{\sinh(m\ell_\alpha/2)}
      \delta_{m\ell_\alpha}
   +\text{a measure supported strictly above }m\ell_\alpha.
 }
 \label{eq:coefffirstatom}
\end{equation}
For each fixed $C$ the gap is locally uniform in $X$, and uniform
on compact subsets of marked Teichm\"uller space. If $C$ has two
or more nonparallel support components, every contributing orbit
has length strictly greater than $\ell_X(C)$; this statement does
not assert that the resulting signed measure is nonzero.
\end{proposition}

\begin{proof}
Lemma~\ref{lem:coefftriangularity} gives the lower support bound.
At $m\ell_\alpha$ the only contributing pair is
$(\delta,k)=(\alpha,m)$, with coefficient $1$. This gives the
stated mass, whose numerator is the \emph{primitive} length.
Every other contributing pair has strictly greater length.
Local finiteness gives a gap at each $X$.

For uniformity on a compact set of marked metrics, use uniform
bilipschitz comparison with one reference metric. The target lengths
are bounded, and only finitely many orbit labels can have lengths
within one unit above them somewhere on the compact set. The
strictly positive differences for those finitely many labels have
a positive minimum; all remaining labels have a gap at least one.
The argument also applies to nonparallel $C$, without requiring a
nonzero first mass after cancellation.
\end{proof}

For $m=1$ the first atom determines both the simple length and its
normalization. Choose a fixed smooth cutoff $\chi$, equal to one
near $\ell_\alpha$ and supported away from every other atom in a
small neighborhood of $X$. Then
\begin{equation}
 \boxed{
 \ell_X(\alpha)
 =\frac{\langle\sigma_\alpha^X,L\chi(L)\rangle}
        {\langle\sigma_\alpha^X,\chi(L)\rangle},\qquad
 \langle\sigma_\alpha^X,\chi\rangle
 =\frac{\ell_\alpha}{\sinh(\ell_\alpha/2)}.
 }
 \label{eq:coeffmomentratio}
\end{equation}
This is an exact localized identity, with no limiting interchange.
Figure~\ref{fig:mono-reconnection-atoms} relates the local reconnections
to this separate first-atom statement for a simple support.

% BEGIN INLINE FIGURE: 05_reconnection_atoms.tex
\begin{figure}[htbp]
\centering
\begin{tikzpicture}[font=\small,>=Latex,line cap=round]
\begin{scope}[shift={(-4.5,1.6)}]
\draw[mpgray!65] (0,0) circle (1.1);
\node[font=\bfseries\small] at (0,1.55) {Crossing};
\draw[mpblue,line width=1.15pt] (0:1.1)--(180:1.1) (60:1.1)--(240:1.1);
\draw[mpochre] (.37,0) arc (0:60:.37);
\node[text=mpochre,font=\scriptsize] at (30:.58) {$\theta$};
\node[text=mpgray] at (0,-1.5) {total length $4r$};
\fill[mpblue] (0:1.1) circle (1.6pt);
\fill[mpblue] (60:1.1) circle (1.6pt);
\fill[mpblue] (180:1.1) circle (1.6pt);
\fill[mpblue] (240:1.1) circle (1.6pt);
\end{scope}
\begin{scope}[shift={(0,1.6)}]
\draw[mpgray!65] (0,0) circle (1.1);
\node[font=\bfseries\small] at (0,1.55) {First reconnection};
\draw[mpteal,line width=1.15pt] (1.10000000,0.00000000) arc[start angle=-132.84678884,delta angle=-34.30642232,radius=1.86486286];
\draw[mpteal,line width=1.15pt] (-1.10000000,0.00000000) arc[start angle=47.15321116,delta angle=-34.30642232,radius=1.86486286];
\node[text=mpteal] at (0,-1.5) {$L_0<4r$};
\fill[mpblue] (0:1.1) circle (1.6pt);
\fill[mpblue] (60:1.1) circle (1.6pt);
\fill[mpblue] (180:1.1) circle (1.6pt);
\fill[mpblue] (240:1.1) circle (1.6pt);
\end{scope}
\begin{scope}[shift={(4.5,1.6)}]
\draw[mpgray!65] (0,0) circle (1.1);
\node[font=\bfseries\small] at (0,1.55) {Second reconnection};
\draw[mpteal,line width=1.15pt] (1.10000000,0.00000000) arc[start angle=107.18035662,delta angle=25.63928676,radius=4.29338334];
\draw[mpteal,line width=1.15pt] (0.55000000,0.95262794) arc[start angle=-47.18035662,delta angle=-25.63928676,radius=4.29338334];
\node[text=mpteal] at (0,-1.5) {$L_\infty<4r$};
\fill[mpblue] (0:1.1) circle (1.6pt);
\fill[mpblue] (60:1.1) circle (1.6pt);
\fill[mpblue] (180:1.1) circle (1.6pt);
\fill[mpblue] (240:1.1) circle (1.6pt);
\end{scope}
  \draw[->,mpgray!80] (-5.75,-2.55)--(5.75,-2.55) node[right] {$\ell$};
  \node[text=mpgray] at (0,-.65) {Coefficient extraction: the first atom of $\sigma_\alpha^X$};
  \draw[mpblue,line width=1.4pt] (-3.55,-2.55)--(-3.55,-1.35);
  \fill[mpblue] (-3.55,-1.35) circle (2pt);
  \node[above=3pt,text=mpblue,font=\scriptsize] at (-3.55,-1.35)
    {$\ell_\alpha/\sinh(\ell_\alpha/2)$};
  \node[below=4pt] at (-3.55,-2.55) {$\ell_\alpha$};
  \foreach \x/\y in {-1.0/-.6,.55/.45,2.0/-.9,3.85/.55} {
    \draw[mpgray,line width=.8pt] (\x,-2.55)--(\x,{-2.55-\y});
    \fill[mpgray] (\x,{-2.55-\y}) circle (1.4pt);
  }
  \draw[<->,mpgray!65] (-3.4,-2.0)--node[above,font=\scriptsize,fill=white,inner sep=1pt]
    {a gap for fixed $X$}(-1.15,-2.0);
  \node[align=center,font=\scriptsize,text=mpgray] at (2.1,-3.45)
    {Later atoms can have either sign; their positions and heights here are schematic.};
\end{tikzpicture}
\caption[Crossing resolution and the first coefficient atom]{Reconnection and coefficient extraction are different operations.
The upper pictures show the two noncrossing pairings in a hyperbolic
metric disk; both shorten the original pair of diameters. Reconnection
changes connectivity and is not a negative direction of the length
Hessian within the original loop space. For a simple primitive support
$\alpha$, the lower schematic records the exact first atom of
$\sigma_\alpha^X$ at $\ell_\alpha$, with the stated mass; all remaining
atoms lie strictly above that length. Their depicted positions and
weights are illustrative, not numerical data.}
\label{fig:mono-reconnection-atoms}
\end{figure}
% END INLINE FIGURE

\subsection{An exponential coefficient bound from bracelet positivity}

The passage from compactly supported tests to heat kernels needs a
growth estimate. A direct crossing count gives an exponential in
$\ell_\gamma^2$ and only a small-time conclusion. Bracelet
positivity gives a stronger estimate.

At skein parameter $+1$, Dylan Thurston proves that every diagram
has a sign-coherent expansion in the bracelet
basis~\cite[Theorem 2]{DThurstonPositive}. A bracelet supported on
a simple curve $\beta$ with winding $n$ is the Chebyshev element
$P_n(z_\beta)$~\cite[Proposition 4.4]{DThurstonPositive}.
Distinct bracelet support curves are disjoint; their windings are
included in the length $\ell_X(D)=\sum_i n_i\ell_X(\beta_i)$.
The canonical twisted Fuchsian evaluation associated to $X$ is
positive on bracelet basis elements:
\[
 B_D(X)=\prod_i2\cosh\bigl(n_i\ell_X(\beta_i)/2\bigr)>0.
\]
For a one-component diagram representing $\gamma$, its evaluation
has absolute value $2\cosh(\ell_X(\gamma)/2)$. Sign coherence
therefore gives
\begin{equation}
 \sum_D |b_D(\gamma)| B_D(X)
 =2\cosh\bigl(\ell_X(\gamma)/2\bigr).
 \label{eq:coeffbraceletpositive}
\end{equation}
A spin structure identifies the $+1$ and $-1$ skein algebras by
signs on diagrams. Changing to~\eqref{eq:coeffbasis} also changes
only basis signs. On a simple generator this replaces $z_\beta$ by
$\epsilon_\beta T_\beta$, with $\epsilon_\beta\in\{1,-1\}$;
the identity $P_n(\epsilon x)=\epsilon^nP_n(x)$ shows that the
bracelet-to-multicurve coefficient magnitudes are unchanged.
All the absolute estimates below are consequently
valid in our unnegated trace basis. The positive Fuchsian evaluation
is an auxiliary algebraic estimate; the Laplacians defining
$\sigma_C^X$ still use unitary twists.

\begin{proposition}[Fixed-coefficient bound]
\label{prop:coeffgrowth}
Write $C=\sum_{j=1}^q r_j\beta_j$ with distinct support curves.
For $r\geq1$ and $\ell>0$ put
\[
 M_r(\ell)=\sup_{n\geq1}
 \frac{|[x^r]P_n(x)|}{2\cosh(n\ell/2)},\qquad
 K_{C,X}=\prod_{j=1}^q M_{r_j}(\ell_X(\beta_j)).
\]
Set $K_{\varnothing,X}=1$. These constants are finite, and
\begin{equation}
 \boxed{
 |a_C(\gamma)|\leq
 K_{C,X}\,2\cosh\bigl(\ell_X(\gamma)/2\bigr)
 }
 \label{eq:coeffexponentialbound}
\end{equation}
for every nontrivial closed curve, including iterates.
\end{proposition}

\begin{proof}
When $n\geq r\geq1$ have the same parity,
\[
 |[x^r]P_n(x)|
 =\frac{2n}{n+r}\binom{(n+r)/2}{r}
 \leq\frac{2n^r}{r!};
\]
otherwise the coefficient is zero. Consequently
\begin{equation}
 M_r(\ell)\leq
 \frac{2}{r!}\left(\frac{2r}{e\ell}\right)^r<\infty.
 \label{eq:coeffMrbound}
\end{equation}
To contribute to $C$, a bracelet element must contain every
$\beta_j$ in its support. Additional support curves must contribute
constant terms, which are either zero or $\pm2$. Their absolute
constant coefficient divided by $2\cosh(n\ell/2)$ is at most one.
It follows that
\[
 |[T_C]B_D|\leq K_{C,X}B_D(X).
\]
Apply this bound to~\eqref{eq:coeffbraceletpositive} and use the
triangle inequality. Spin and basis signs do not affect the estimate.
\end{proof}

For a single simple $\alpha$ one can use
\[
 K_{\alpha,X}
 =\sup_{\substack{n\geq1\\n\text{ odd}}}
   \frac{n}{2\cosh(n\ell_\alpha/2)}
 \leq\frac{2}{e\ell_\alpha}.
\]
The estimates are uniform for a fixed target $C$ on compact subsets
of Teichm\"uller space. Related positivity and Chebyshev estimates
for analytic lamination expansions on punctured planar surfaces
appear in Basok--Chelkak~\cite[Section 4.4, Theorems 4.7 and 4.9]{BasokChelkak}. Their
full analytic expansion theorem is not needed here.

\begin{corollary}[Growth of the signed measure]
\label{cor:coeffmeasuregrowth}
For each fixed $X,C$ there is a constant $B_{C,X}$ such that
\begin{equation}
 |\sigma_C^X|([0,R])\leq B_{C,X}(1+R)e^R.
 \label{eq:coeffmeasuregrowth}
\end{equation}
\end{corollary}

\begin{proof}
Put $\varepsilon=\operatorname{sys}(X)>0$ and
$L=k\ell_\delta$. Each orbital summand has absolute mass at most
\[
 \frac{\ell_\delta}{\sinh(L/2)}
 K_{C,X}2\cosh(L/2)
 \leq2K_{C,X}\coth(\varepsilon/2)L.
\]
The number of primitive/iterate classes of length at most $R$ is
$O_X(e^R)$. One elementary proof conjugates each class so its axis
passes a bounded distance from a fixed point in $\HH^2$; the chosen
element then moves that point by at most $R+O_X(1)$. Orbit points
are uniformly separated, and hyperbolic disk area is $O(e^R)$.
Summing the displayed mass bound proves the claim. No prime-geodesic
asymptotic is required.
\end{proof}

\subsection{Convergent heat and Laplace transforms}

Define
\begin{align}
 G_C^X(t)&=\int_0^\infty e^{-L^2/(4t)}\,\dd\sigma_C^X(L),
 &H_C^X(t)&=\frac{e^{-t/4}}{\sqrt{4\pi t}}G_C^X(t),
 \label{eq:coeffheatdefinition}\\
 D_C^X(s)&=\int_0^\infty e^{-sL}\,\dd\sigma_C^X(L).
 &&\label{eq:coeffLaplacedefinition}
\end{align}
By~\eqref{eq:coeffmeasuregrowth}, $G_C$ and $H_C$ are absolutely
convergent for every $t>0$, and holomorphic for $\Re t>0$ with
the principal square root. The Laplace transform is absolutely
convergent and holomorphic for $\Re s>1$. Differentiation in these
variables is permitted locally uniformly in the stated domains.

There is also a convergent expansion of the whole heat trace.
For a finite character polynomial set
\[
 \|f\|_2=\sum_C|[T_C]f|\,2^{\#C},\qquad
 c_X=\frac{\log(1+\sqrt2)}{\operatorname{sys}(X)}.
\]
The subscript $2$ denotes the trace bound $|\tr\rho(\alpha)|\leq2$;
this is a weighted coefficient norm, not an $L^2$ norm.
The Chebyshev recurrence gives
\[
 \sum_r|[x^r]P_n(x)|2^r\leq2(1+\sqrt2)^n.
\]
A bracelet has at most $3g-3$ distinct support curves, so
$\|B_D\|_2\leq2^{3g-3}e^{c_X\ell_X(D)}$.
In~\eqref{eq:coeffbraceletpositive} the nonzero coefficients are
integers, and $B_D(X)\geq e^{\ell_X(D)/2}$. Hence
$\ell_X(D)\leq\ell_X(\gamma)+2\log2$ for every occurring $D$,
and $\sum_D|b_D|\leq2e^{\ell_X(\gamma)/2}$. Therefore
\begin{equation}
 \|\tr_\gamma\|_2\leq
 C_X e^{(1/2+c_X)\ell_X(\gamma)},\qquad
 C_X=2^{3g-2}e^{2c_X\log2}.
 \label{eq:coeffallchannelbound}
\end{equation}
Gaussian decay dominates this exponential bound and the orbit count.
It follows that
\begin{equation}
 \boxed{
 \operatorname{Tr}e^{-t\Delta_\rho}-2I_X(t)
 =\sum_{C\in\mathscr C(S)}H_C^X(t)T_C(\rho),\qquad t>0,
 }
 \label{eq:coefffullheatexpansion}
\end{equation}
where
\[
 I_X(t)=\frac{\operatorname{Area}(X)}{4\pi}
  \int_\RR r\tanh(\pi r)e^{-t(r^2+1/4)}\,\dd r,
 \qquad \sum_C|H_C^X(t)|2^{\#C}<\infty.
\]
The series is uniform on the unitary character variety and locally
uniform for $\Re t>0$. It is obtained by taking the limit of the
finite length-cutoff expansions; the preceding absolute estimate
justifies the rearrangement. For a compact set of such $t$,
$\Re(1/t)$ has a positive lower bound, so the same Gaussian
majorant proves local uniformity and termwise holomorphy. On the
spectral side, smooth cutoffs of the Gaussian converge in the
Schwartz topology on the real spectral axis and at the finitely
many imaginary spectral parameters; the compact-surface Weyl bound
justifies passing to the heat trace for each unitary twist.

This construction does not assert uniqueness among arbitrary
infinite expansions of a continuous function. For example,
$\|\tr\rho(\alpha^{2n+1})\|_\infty\leq2$, while its
$T_\alpha$ coefficient is $(-1)^n(2n+1)$. Coefficient extraction
is therefore not bounded for the uniform norm on unitary characters.
Our coefficients are fixed by the finite spectral tests first;
the estimates then justify their heat transform.

\begin{proposition}[Simple-curve semiclassical term]
\label{prop:coeffheatfirst}
Let $\alpha$ be simple and $\ell=\ell_X(\alpha)$. Choose
$b>\ell$ so that every other orbital contribution to
$\sigma_\alpha^X$ has length at least $b$. Then
\begin{align}
 G_\alpha^X(t)
 &=\frac{\ell}{\sinh(\ell/2)}e^{-\ell^2/(4t)}
   +O\bigl(e^{-b^2/(4t)}\bigr),\notag\\
 H_\alpha^X(t)
 &=\frac{e^{-t/4}}{\sqrt{4\pi t}}
   \frac{\ell}{\sinh(\ell/2)}e^{-\ell^2/(4t)}
   \left[1+O\!\left(e^{-(b^2-\ell^2)/(4t)}\right)\right],
 \label{eq:coeffheatfirst}\\
 D_\alpha^X(s)
 &=\frac{\ell}{\sinh(\ell/2)}e^{-s\ell}
   \left[1+O\bigl(e^{-s(b-\ell)}\bigr)\right]
   \quad(s\to+\infty).
 \label{eq:coeffLaplacefirst}
\end{align}
In particular,
\begin{equation}
 \boxed{
 \ell^2=-\lim_{t\downarrow0}4t\log G_\alpha^X(t),\qquad
 \ell=-\lim_{s\to+\infty}s^{-1}\log D_\alpha^X(s).
 }
 \label{eq:coefflengthlimits}
\end{equation}
Both logarithms are real for sufficiently small $t$, respectively
sufficiently large positive $s$.
\end{proposition}

\begin{proof}
Separate the first atom in~\eqref{eq:coefffirstatom}. For its
remaining measure $\tau$, supported on $[b,\infty)$, fix $t_0>0$.
For $0<t\leq t_0$,
\[
 \int e^{-L^2/(4t)}\,\dd|\tau|(L)
 \leq e^{-b^2/(4t)+b^2/(4t_0)}
       \int e^{-L^2/(4t_0)}\,\dd|\tau|(L).
\]
The last integral is finite. The same argument with a fixed real
$s_0>1$ gives the Laplace remainder. The leading masses are
positive, so the leading terms eventually dominate the signed tails.
\end{proof}

After removing the explicit $e^{-t/4}$ factor, the correction to
the single-geodesic term is exponentially small, rather than a new
power series in $t$. Thus the constant-curvature heat coefficient
has an exact isolated-orbit prefactor, followed by contributions at
larger length actions. For $m\alpha$ replace the action by
$(m\ell_\alpha)^2/(4t)$ and the mass by
$\ell_\alpha/\sinh(m\ell_\alpha/2)$.

\subsection{All iterates of a simple curve}

Write $\sigma_C^{X,[\delta]}$ for the contribution of one primitive
unoriented class $\delta$ and all of its iterates. The polynomial
identity
\begin{equation}
 P_k(x)=\sum_{j=0}^{\lfloor k/2\rfloor}
 (-1)^j\frac{k(k-j-1)!}{j!(k-2j)!}x^{k-2j},\qquad k\geq1,
 \label{eq:coeffChebyshev}
\end{equation}
follows from the recurrence. If $\alpha$ is simple, the coefficient
of $m\alpha$ vanishes unless $k=m+2j$, and for $m\geq1$ is
\begin{equation}
 b_{m+2j,m}=(-1)^j
  \frac{(m+2j)(m+j-1)!}{j!\,m!}.
 \label{eq:coeffsimplepowers}
\end{equation}
In particular, with $\ell=\ell_\alpha$,
\begin{equation}
 \boxed{
 \sigma_\alpha^{X,[\alpha]}
 =\ell\sum_{j\geq0}
   \frac{(-1)^j(2j+1)}{\sinh((2j+1)\ell/2)}
   \delta_{(2j+1)\ell}.
 }
 \label{eq:coeffoddpowers}
\end{equation}
The first masses have multipliers $1,-3,5,-7,\ldots$ at
$\ell,3\ell,5\ell,7\ell,\ldots$. For $2\alpha$ the multipliers
are $1,-4,9,-16,\ldots$ at $2\ell,4\ell,6\ell,8\ell,\ldots$.
These are explicit later atoms, not changes to the first-atom mass.

For $m\geq1$ the coefficient generating function is
\begin{equation}
 B_m(q):=\sum_{k\geq1}b_{k,m}q^k
 =\frac{q^m(1-q^2)}{(1+q^2)^{m+1}}.
 \label{eq:coeffpowergenerating}
\end{equation}
Indeed,
$\sum_{k\geq1}P_k(x)q^k=(xq-2q^2)/(1-xq+q^2)$,
and extracting $x^m$ gives this formula. Expanding the hyperbolic
denominator yields, for $\Re s>-1/2$,
\begin{align}
 D_{m\alpha}^{X,[\alpha]}(s)
 &=2\ell\sum_{n\geq0}B_m(q_n),
 &q_n&=e^{-(s+n+1/2)\ell},\notag\\
 &=2^{1-m}\ell\sum_{n\geq0}
   \sech^m(u_n)\tanh(u_n),
 &u_n&=(s+n+1/2)\ell.
 \label{eq:coeffpowerLaplace}
\end{align}
This larger convergence domain concerns the one fixed primitive
family; the proved half-plane for the full $D_C$ is $\Re s>1$.

The heat contribution of the same family is
\begin{align}
 H_\alpha^{X,[\alpha]}(t)
 ={}&\frac{e^{-t/4}}{\sqrt{4\pi t}}
 \frac{\ell e^{-\ell^2/(4t)}}{\sinh(\ell/2)}\notag\\
 &\mathrel{}\times\left[
 1-3\frac{\sinh(\ell/2)}{\sinh(3\ell/2)}e^{-2\ell^2/t}
 +5\frac{\sinh(\ell/2)}{\sinh(5\ell/2)}e^{-6\ell^2/t}
 -\cdots\right].
 \label{eq:coeffpowerheat}
\end{align}
Other primitive classes may supply corrections before these particular
iterate actions. Once $\ell_\alpha$ is recovered, subtracting the
$j\geq1$ terms of~\eqref{eq:coeffoddpowers} removes exactly the
nonprimitive powers of the simple curve $\alpha$; it leaves powers
of nonsimple curves that resolve with an $\alpha$ coefficient.

\subsection{A figure-eight family and an exact cancellation}

Take an embedded geodesic pair of pants with three distinct essential
cuffs $a,b,c$, and $c=(ab)^{-1}$. Set
\[
 x=T_a,\qquad y=T_b,\qquad z=T_c,\qquad
 \gamma=a^{-1}b,
 \qquad T_\gamma=xy-z.
\]
These are functions of arbitrary unitary holonomy. The minus sign
therefore differs from the positive hyperbolic length convention
in~\eqref{eq:figuretrace}. Powers are given exactly by
$T_{\gamma^k}=P_k(xy-z)$. For example,
\begin{align*}
 T_\gamma&=xy-z,\\
 T_{\gamma^2}&=x^2y^2-2xyz+z^2-2,\\
 T_{\gamma^3}&=x^3y^3-3x^2y^2z+3xyz^2-z^3-3xy+3z.
\end{align*}
Only multicurves $C=ra+rb+mc$ occur. Put $d=r+m\geq1$.
Their coefficients vanish unless $k=d+2j$, and in that case
\begin{equation}
 \boxed{
 a_{ra+rb+mc}(\gamma^k)
 =(-1)^{j+m}\frac{k(k-j-1)!}{j!\,r!\,m!}
 =(-1)^m\binom{r+m}{r}b_{k,r+m}.
 }
 \label{eq:coefffigurepowers}
\end{equation}
This follows by substituting $xy-z$ into~\eqref{eq:coeffChebyshev}
and expanding the relevant monomial. Consequently
\[
 \sum_{k\geq1}a_{ra+rb+mc}(\gamma^k)q^k
 =(-1)^m\binom{r+m}{r}
 \frac{q^{r+m}(1-q^2)}{(1+q^2)^{r+m+1}}.
\]

For $r+m=1$ the two target coefficients are opposite for every
iterate. If $L=\ell_\gamma$, then
\begin{equation}
 \boxed{
 \sigma_{a+b}^{X,[\gamma]}
 =-\sigma_c^{X,[\gamma]}
 =L\sum_{j\geq0}
 \frac{(-1)^j(2j+1)}{\sinh((2j+1)L/2)}
 \delta_{(2j+1)L}.
 }
 \label{eq:coefffigurecancellation}
\end{equation}
Thus $\sigma_c^X+\sigma_{a+b}^X$ cancels the entire chosen
figure-eight family, including all its powers. This does not assert
that the full sum vanishes: other primitive classes contribute to
both coefficients. It gives an exact family-by-family cancellation
which is preserved by either transform in
\eqref{eq:coeffheatdefinition}--\eqref{eq:coeffLaplacedefinition}.

The distinct-cuff assumption makes the monomials above distinct
multicurve basis elements. If cuff images coincide, as in some
one-holed-torus configurations of Section~\ref{sec:figures}, the
corresponding variables must first be identified and equal basis
terms collected; the displayed coefficient indexing cannot be
carried over unchanged.

\subsection{Metric variation and the twist derivative}

For a fixed topological $C$, the map
$X\mapsto\sigma_C^X$ is smooth as a distribution-valued map.
After applying any compactly supported test function, it is locally
a finite sum of smooth length functions. It is not asserted to be
real analytic for arbitrary smooth test functions.

For an orbital label $(\delta,k)$ write
\[
 L=k\ell_\delta,\qquad
 w=\ell_\delta\operatorname{csch}(L/2).
\]
For a tangent vector $v$ to Teichm\"uller space one has
\begin{equation}
 D_v\sigma_C^X
 =\sum_{\delta,k}a_C(\delta^k)
  \left[(D_vw)\delta_L-w(D_vL)\delta'_L\right],
 \label{eq:coeffmetricvariation}
\end{equation}
where $\langle\delta'_L,\varphi\rangle=-\varphi'(L)$ and
\[
 D_vL=kD_v\ell_\delta,\qquad
 D_vw=\operatorname{csch}(L/2)
       \left(1-\frac L2\coth\frac L2\right)D_v\ell_\delta.
\]
The coefficients $a_C$ are constant under metric variation. This
formula describes movement of the atoms as well as variation of
their masses; omitting the $\delta'_L$ term would lose the length
variation itself.

For a twist of displacement $\tau_\beta$ about a simple geodesic
$\beta$, Wolpert's cosine formula gives
\begin{equation}
 \frac{\partial\ell_\delta}{\partial\tau_\beta}
 =\sum_{p\in\delta\cap\beta}\cos\theta_p,
 \label{eq:coeffWolpert}
\end{equation}
with the consistent right-twist angle convention and intersections
counted with multiplicity~\cite{WolpertTwist}. Substitution in
\eqref{eq:coeffmetricvariation} gives an exact variation formula
for every coefficient distribution. For a simple $\alpha$, one
can equivalently differentiate the localized ratio
\eqref{eq:coeffmomentratio}; it recovers~\eqref{eq:coeffWolpert}
from two spectral coefficient moments, without an infinite limit.
The first-variation argument also applies when $\delta$ is immersed:
cut it at its crossings with $\beta$. Interior variations of the
geodesic arcs cancel by stationarity; the relative tangential
displacement at each paired endpoint contributes the projection
$\cos\theta_p$. Summing over the passages through the cut gives
the formula with precisely the stated multiplicities. If $\delta$
is a power of $\beta$, its twist derivative is zero and the sum is
interpreted as empty.

The compact-set gap in Proposition~\ref{prop:coefffirstatom} also
permits differentiating the leading heat asymptotic locally.
The leading logarithmic response is
\begin{equation}
 D_v\log H_\alpha^X(t)
 =\left[-\frac{\ell_\alpha}{2t}
          +\frac1{\ell_\alpha}
          -\frac12\coth\frac{\ell_\alpha}{2}\right]
    D_v\ell_\alpha
  +O\bigl(t^{-1}e^{-c/t}\bigr)
 \label{eq:coeffheatvariation}
\end{equation}
for some local $c>0$. To justify the differentiated tail, metric
length derivatives on a compact parameter neighborhood are bounded
by constants times the length, while
\eqref{eq:coeffexponentialbound} is locally uniform. Gaussian
decay absorbs the additional polynomial factors.

\subsection{The logarithmic product and primitive extraction}

A useful connection to the Selberg product keeps the orientation
normalization explicit. Define, for $\Re v>1$,
\begin{equation}
 \mathcal Z_\rho^X(v)
 =\prod_{\delta\in\mathcal P_{\rm un}}\prod_{n\geq0}
 \det\left(I-e^{-(v+n)\ell_\delta}\rho(\delta)\right).
 \label{eq:coeffunorientedproduct}
\end{equation}
This definition uses unoriented primitives; a product over both
oriented conjugacy classes squares it. Expanding the logarithm gives
\begin{equation}
 2\left.\frac{\partial}{\partial v}
     \log\mathcal Z_\rho^X(v)\right|_{v=s+1/2}
 =\sum_{\delta,k}
   \frac{\ell_\delta e^{-sk\ell_\delta}}
        {\sinh(k\ell_\delta/2)}\tr\rho(\delta^k).
 \label{eq:coeffproductderivative}
\end{equation}
With coefficients defined by finite length cutoffs, the $T_C$
coefficient on the right is $D_C^X(s)$ for $\Re s>1$.
The coefficient estimate proves this assertion; it is not an
appeal to continuity of coefficient extraction in the uniform norm.
No meromorphic continuation of the individual $D_C$ is established
by this calculation.

Removing all iterates by a global Chebyshev substitution is not
legitimate in the rank-two character ring. If
$x=\tr A$, $y=\tr B$, $u=\tr(AB)$, and $v=\tr(AB^{-1})$,
then $u+v=xy$, but
\begin{equation}
 P_2(u)+P_2(v)-P_2(x)P_2(y)
 =2\bigl(\tr[A,B]-2\bigr),
 \label{eq:coeffAdamsobstruction}
\end{equation}
which is generally nonzero. Thus sending every trace to its
Chebyshev polynomial does not define a multiplicative operation
on this character ring.

There is an exact repair if extra twisting data are supplied.
For the original $\mathrm{SU}(2)$ representation $\rho$, work in
the virtual representation ring generated by its symmetric powers.
Let $\psi^k\rho$ denote its Adams operation, characterized by
$\chi_{\psi^k\rho}(\gamma)=\chi_\rho(\gamma^k)$. For the
standard rank-two representation,
\[
 \psi^k\rho=\operatorname{Sym}^k\rho-
             \operatorname{Sym}^{k-2}\rho\quad(k\geq2),
 \qquad \psi^1\rho=\rho.
\]
Trace formulas extend additively to virtual representations.
All twists used here have inverse-invariant characters, so the
unoriented normalization is preserved. Let $\Sigma_\rho$ denote the character-valued geometric measure
on the right of~\eqref{eq:coefftwistedSelberg}, and set
\[
 Q_\rho=\sinh(L/2)\Sigma_\rho,\qquad
 P_\rho=\sum_{\delta\in\mathcal P_{\rm un}}
         \ell_\delta\chi_\rho(\delta)\delta_{\ell_\delta}.
\]
For $D_k\delta_L=\delta_{kL}$,
\[
 Q_\rho=\sum_{k\geq1}D_kP_{\psi^k\rho},\qquad
 \boxed{P_\rho=\sum_{k\geq1}\mu(k)D_kQ_{\psi^k\rho}.}
\]
These sums are locally finite, and the second identity follows from
$\psi^j\psi^k=\psi^{jk}$ and divisor M\"obius inversion. It requires
spectra for higher symmetric-power twists, which are additional
data beyond the original rank-two family. The explicit subtraction
in~\eqref{eq:coeffoddpowers} needs no such enlargement.

\subsection{Scope and further questions}

The construction now supplies signed coefficient distributions,
an isolated first simple-curve atom, globally convergent heat
transforms, a Laplace half-plane, explicit iterate families, a
figure-eight cancellation, and metric derivatives. It has not
produced a positive measure with the weights of simple-geodesic
counting, nor a pointwise constant McShane series.

The scalar limitation is genuine: length-isospectral closed
hyperbolic surfaces can have different simple length spectra
\cite{AougabLahnLovingMiller}. The labelled holonomy family retains
the information used by~\eqref{eq:coeffmomentratio}.

For a counting comparison, let $d=6g-6$ and normalize Thurston
measure by the integral multicurve lattice. If
\[
 B(X)=\mu_{\rm Th}\{\lambda:\ell_X(\lambda)\leq1\},
\]
ordinary lattice counting gives
$\#\{C\in\mathcal{ML}_{\ZZ}:\ell_X(C)\leq R\}\sim B(X)R^d$
\cite[Proposition 3.1]{MirzSimple}. Its Laplace consequence is
\[
 \sum_{C\in\mathcal{ML}_{\ZZ}}e^{-t\ell_X(C)}
 \sim\Gamma(d+1)B(X)t^{-d}\qquad(t\downarrow0).
\]
This is a target for a measure comparison, not a consequence of
the coefficient heat expansion. It counts all weighted multicurves;
connected or fixed-orbit counts require additional information.
The large-length limit here is also different from the
minimum-selection limits in~\eqref{eq:coefflengthlimits}.

Two concrete further problems are to continue $D_C(s)$ beyond its
proved half-plane, with control of the coefficient operation, and
to find combinations of coefficients whose cancellations organize
globally over pants. Formula~\eqref{eq:coefffigurecancellation}
provides an exact local model for the second problem. An identity
constant in $X$ would require an additional normalization or
geometric partition; the existence of these transforms alone does
not supply that constancy. Section~\ref{sec:quantitativecounting}
strengthens the lattice benchmark and gives explicit orbital gaps,
uniform recovery estimates, and a conditional counting criterion.


\clearpage
